\section*{Chapter \ref{ch:b1:rock-forming-minerals} --- Rock-Forming Minerals}

\begin{solution}{exo:b1:rock-forming-minerals:1}
\ce{SiO4^{4-}}: \omterm{def:b1:rock-forming-minerals:silicate-classes}{nesosilicate}, $s = 0$. \ce{Si2O7^{6-}}: \omterm{def:b1:rock-forming-minerals:silicate-classes}{sorosilicate},
$s = 1$. \ce{Si6O18^{12-}}: \omterm{def:b1:rock-forming-minerals:silicate-classes}{cyclosilicate}, $s = 2$. \ce{Si2O5^{2-}}:
\omterm{def:b1:rock-forming-minerals:silicate-classes}{phyllosilicate}, $s = 3$. \ce{SiO2}: \omterm{def:b1:rock-forming-minerals:silicate-classes}{tectosilicate}, $s = 4$. \ce{SiO3^{2-}}
has $4 - s/2 = 3$, so $s = 2$: a ring (\omterm{def:b1:rock-forming-minerals:silicate-classes}{cyclosilicate}) or a single chain
(\omterm{def:b1:rock-forming-minerals:silicate-classes}{inosilicate}); the ratio alone cannot tell them apart.
\end{solution}

\begin{solution}{exo:b1:rock-forming-minerals:2}
Forsterite: $2\times24.305 + 28.085 + 4\times15.999 =
\qty{140.691}{g.mol^{-1}}$; fayalite: $2\times55.845 + 28.085 +
4\times15.999 = \qty{203.771}{g.mol^{-1}}$. Molar volumes:
$140.691/3.275 = \qty{42.96}{cm^3.mol^{-1}}$ and $203.771/4.392 =
\qty{46.40}{cm^3.mol^{-1}}$.
\end{solution}

\begin{solution}{exo:b1:rock-forming-minerals:3}
$(0.78 - 0.72)/0.72 = +\qty{8.3}{\percent}$ for \ce{Fe^{2+}};
$-\qty{4.2}{\percent}$ for \ce{Ni^{2+}}; $+\qty{15.3}{\percent}$ for
\ce{Mn^{2+}}; $+\qty{38.9}{\percent}$ for \ce{Ca^{2+}}. Iron and nickel,
under \qty{15}{\percent}, replace magnesium freely; manganese, just above, is
limited; calcium, beyond \qty{30}{\percent}, hardly enters. All four have the
same charge as magnesium.
\end{solution}

\begin{solution}{exo:b1:rock-forming-minerals:4}
$(3.51 - 2.16)/2.16 = \qty{62.5}{\percent}$ from graphite to diamond;
$(2.93 - 2.71)/2.71 = \qty{8.1}{\percent}$ from calcite to aragonite. The
denser \omterm{def:b1:rock-forming-minerals:polymorph}{polymorphs}, diamond and aragonite, are the high-pressure forms.
\end{solution}

\begin{solution}{exo:b1:rock-forming-minerals:5}
With $x = 0.30$ in $\mathrm{Na}_{1-x}\mathrm{Ca}_x\mathrm{Al}_{1+x}\mathrm{Si}_{3-x}\mathrm{O}_8$:
$\mathrm{(Na_{0.7}Ca_{0.3})(Al_{1.3}Si_{2.7})O_8}$. Charges:
$0.7 + 2\times0.3 + 3\times1.3 + 4\times2.7 = 0.7 + 0.6 + 3.9 + 10.8 = +16$,
and $8\times(-2) = -16$.
\end{solution}

\begin{solution}{exo:b1:rock-forming-minerals:6}
A repeating group has two tetrahedra sharing three corners (those whose apex
joins the two chains) and two sharing two. The four tetrahedra have
$4\times4 = 16$ corners, of which $3 + 3 + 2 + 2 = 10$ are shared, that is
$10/2 = 5$ shared oxygens; the $16 - 10 = 6$ others are unshared. Total
$5 + 6 = 11$ oxygens for 4 silicons: \ce{Si4O11}, with charge
$4\times4 - 11\times2 = -6$. The proposition gives the same with
$s = 5/2$: $4 - 5/4 = 11/4$ oxygens per silicon, charge $5/2 - 4 = -3/2$ per
silicon.
\end{solution}

\begin{solution}{exo:b1:rock-forming-minerals:7}
Moles in \qty{100}{g}: \ce{SiO2} $40.20/60.083 = 0.66907$; \ce{MgO}
$47.10/40.304 = 1.16862$; \ce{FeO} $12.20/71.844 = 0.16981$; \ce{NiO}
$0.40/74.692 = 0.00536$. Oxygens: $2\times0.66907 + 1.16862 + 0.16981 +
0.00536 = 2.68194$; $f = 4/2.68194 = 1.49146$. Per four oxygens: Si
\num{0.998}, Mg \num{1.743}, Fe \num{0.253}, Ni \num{0.008}. The
octahedral sites hold $1.743 + 0.253 + 0.008 = 2.004$, the tetrahedral site
\num{0.998}: a good analysis, \ce{(Mg_{1.743}Fe_{0.253}Ni_{0.008})SiO4}.
$\mathrm{Fo} = 100\times1.743/(1.743 + 0.253) = 87.3$. The factor $f$
multiplies Mg and Fe alike, so it cancels in the ratio:
$\mathrm{Fo} = 100\,(47.10/40.304)/(47.10/40.304 + 12.20/71.844) =
100\times1.16862/1.33843 = 87.3$.
\end{solution}

\begin{solution}{exo:b1:rock-forming-minerals:8}
Positive: $1 + 2\times3 + 3\times4 + 3 = +22$; negative:
$10\times(-2) + 2\times(-1) = -22$. A sheet of composition \ce{Si4O10} would
carry $-4$ (\cref{prop:b1:rock-forming-minerals:si-o}, $s = 3$); replacing one
silicon in four by aluminium adds one negative charge, so \ce{(Si3Al)O10}
carries $-5$. The two octahedral aluminiums with the two hydroxyls,
$6 - 2 = +4$, balance four of the five; the potassium ion between the sheets
balances the fifth.
\end{solution}

\begin{solution}{exo:b1:rock-forming-minerals:9}
(a) $M = (140.691 + 203.771)/2 = \qty{172.231}{g.mol^{-1}}$, $V = (42.96 +
46.40)/2 = \qty{44.68}{cm^3.mol^{-1}}$, $\rho = 172.231/44.68 =
\qty{3.855}{g.cm^{-3}}$; the mean of the densities is $(3.275 + 4.392)/2 =
\qty{3.834}{g.cm^{-3}}$, \qty{0.021}{g.cm^{-3}} less. (b) Solve
$203.771 + x(140.691 - 203.771) = 3.600\,\bigl(46.40 + x(42.96 - 46.40)\bigr)$:
$203.771 - 63.080x = 167.040 - 12.384x$, so $x = 36.731/50.696 = 0.7245$: the
\omterm{def:b1:rock-forming-minerals:groups}{olivine} is $\text{Fo}_{72}$ (to the nearest per cent).
\end{solution}

\begin{solution}{exo:b1:rock-forming-minerals:10}
With $M_i = \rho_i V_i$,
\[
  \rho(x) - \bigl(x\rho_1 + (1 - x)\rho_2\bigr)
  = \frac{x\rho_1V_1 + (1 - x)\rho_2V_2 - \bigl(x\rho_1 + (1 - x)\rho_2\bigr)\bigl(xV_1 + (1 - x)V_2\bigr)}{V(x)} .
\]
Expanding the product, the numerator is $x\rho_1V_1 - x^2\rho_1V_1 +
(1 - x)\rho_2V_2 - (1 - x)^2\rho_2V_2 - x(1 - x)(\rho_1V_2 + \rho_2V_1)$, and
$x - x^2 = x(1 - x)$, $(1 - x) - (1 - x)^2 = x(1 - x)$, so it is
$x(1 - x)(\rho_1V_1 + \rho_2V_2 - \rho_1V_2 - \rho_2V_1) =
x(1 - x)(\rho_1 - \rho_2)(V_1 - V_2)$. For \omterm{def:b1:rock-forming-minerals:groups}{olivine} (1 = forsterite) both
differences are negative, $3.275 - 4.392 = -1.117$ and $42.96 - 46.40 =
-3.44$, so their product is positive and the density lies above the straight
line. At $\text{Fo}_{50}$: $0.25\times1.117\times3.44/44.68 =
\qty{0.0215}{g.cm^{-3}}$, which agrees with \cref{exo:b1:rock-forming-minerals:9} to the rounding of the
molar volumes.
\end{solution}

\begin{solution}{exo:b1:rock-forming-minerals:11}
Total \qty{99.60}{\percent}. Moles of oxide: \ce{SiO2} $0.90042$,
\ce{Al2O3} $0.28344$, \ce{CaO} $0.19973$, \ce{Na2O} $0.08229$,
\ce{K2O} $0.00318$. Oxygens: $1.80084 + 0.85033 + 0.19973 + 0.08229 +
0.00318 = 2.93636$; $f = 8/2.93636 = 2.72446$. Cations (twice the moles for
\ce{Al2O3}, \ce{Na2O}, \ce{K2O}): Si \num{2.453}, Al \num{1.544}, Ca
\num{0.544}, Na \num{0.448}, K \num{0.017}. Tetrahedral sites:
$2.453 + 1.544 = 3.997$; large site: $0.544 + 0.448 + 0.017 = 1.009$. Then
$\mathrm{Al} - 1 = 0.544$, equal to Ca within the rounding: each calcium came
in with one aluminium. $\mathrm{An} = 100\times0.544/1.009 = 53.9$: a
labradorite, $\text{An}_{54}$.
\end{solution}

\begin{solution}{exo:b1:rock-forming-minerals:12}
(a) $M(\ce{Fe3O4}) = 3\times55.845 + 4\times15.999 = \qty{231.531}{g.mol^{-1}}$.
Its three irons, reported as \ce{FeO}, weigh $3\times71.844 = \qty{215.532}{g}$,
that is $215.532/231.531 = \qty{93.09}{\percent}$; the oxygen bound to the
\ce{Fe^{3+}}, beyond what \ce{FeO} accounts for, is missing from the total.
(b) $93.09/71.844 = 1.29572$ mole of \ce{FeO} and of oxygen; $f =
4/1.29572 = 3.08708$; four irons per four oxygens, \ce{Fe4O4}. Magnetite has
three cations for four oxygens. The oxygens were undercounted: each
\ce{Fe^{3+}} holds one and a half oxygens, not one as \ce{FeO} assumes. (c) On three cations, Fe $= 3$
with $a$ \ce{Fe^{2+}} and $b$ \ce{Fe^{3+}}: $a + b = 3$ and $2a + 3b = 8$,
so $b = 2$, $a = 1$: $\mathrm{Fe^{2+}Fe^{3+}_2O_4}$. (d) One \ce{FeO} and
one \ce{Fe2O3} per formula unit: $71.844/231.531 = \qty{31.03}{\percent}$
\ce{FeO} and $159.687/231.531 = \qty{68.97}{\percent}$ \ce{Fe2O3}, total
\qty{100.00}{\percent}.
\end{solution}

\begin{solution}{pb:b1:rock-forming-minerals:1}
\textbf{1.} \ce{SiO2} $28.085 + 2\times15.999 = 60.083$; \ce{FeO} 71.844;
\ce{MnO} 70.937; \ce{MgO} 40.304; \ce{NiO} 74.692; \ce{CaO} 56.077
(\si{g.mol^{-1}}).

\textbf{2.} $39.50 + 16.53 + 0.23 + 43.32 + 0.20 + 0.22 =
\qty{100.00}{\percent}$. A total of \qty{97}{\percent} would have suggested a
component not analysed (water, or oxygen bound to \ce{Fe^{3+}} reported as
\ce{FeO}), or a poor analysis.

\textbf{3.} \ce{SiO2} $39.50/60.083 = 0.65742$; \ce{FeO} 0.23008; \ce{MnO}
0.00324; \ce{MgO} 1.07483; \ce{NiO} 0.00268; \ce{CaO} 0.00392.

\textbf{4.} $2\times0.65742 = 1.31485$ for \ce{SiO2}, the same numbers as in
question 3 for the others; sum $1.31485 + 0.23008 + 0.00324 + 1.07483 +
0.00268 + 0.00392 = 2.62960$.

\textbf{5.} Each oxide has one cation: the numbers of question 3.

\textbf{6.} $f = 4/2.62960 = 1.52114$.

\textbf{7.} Si $0.65742\times1.52114 = 1.000$; Fe 0.350; Mn 0.005; Mg 1.635;
Ni 0.004; Ca 0.006.

\textbf{8.} $1.635 + 0.350 + 0.005 + 0.004 + 0.006 = 2.000$ for the
octahedral sites and 1.000 for the tetrahedral site: exactly the ideal 2 and
1.

\textbf{9.} \ce{(Mg_{1.635}Fe_{0.350}Mn_{0.005}Ni_{0.004}Ca_{0.006})SiO4}.

\textbf{10.} $2\times2.000 + 4\times1.000 = +8.000$ against $4\times(-2) =
-8$. Each cation was counted with the charge of its oxide, so the positive
charges are twice the oxygens by construction: the check cannot fail. The
site totals of question 8 could have.

\textbf{11.} Nickel differs from magnesium by
$(0.69 - 0.72)/0.72 = -\qty{4}{\percent}$, well within \qty{15}{\percent}:
it substitutes freely. Calcium is $(1.00 - 0.72)/0.72 = \qty{39}{\percent}$
larger, beyond \qty{30}{\percent}: it hardly enters.

\textbf{12.} Iron counted as \ce{FeO} when part of it is \ce{Fe^{3+}} gives
too few oxygens per cation, hence more than three cations on four oxygens
(\cref{exo:b1:rock-forming-minerals:12}). The sum is $1.000 + 2.000 = 3.000$: no significant
\ce{Fe^{3+}}.

\textbf{13.} $\mathrm{Fo} = 100\times1.635/(1.635 + 0.350) = 163.5/1.985 =
82.4$.

\textbf{14.} $\mathrm{Mg\#} = 82.4$: all the iron is \ce{Fe^{2+}}
(question 12), so $\mathrm{Fe^{2+}} = \mathrm{Fe}$ and the two definitions
coincide.

\textbf{15.} $\text{Fo}_{82.4}\text{Fa}_{17.6}$, $x = 0.824$; iron fills
\qty{17.6}{\percent} of the octahedral sites in the binary model
($0.350/2.000 = \qty{17.5}{\percent}$ with the minor elements).

\textbf{16.} The factor multiplies Mg and Fe alike and cancels:
$\mathrm{Fo} = 100\times1.07483/(1.07483 + 0.23008) = 82.4$.

\textbf{17.} $40.5/40.304 = 1.00486$ and $20.0/71.844 = 0.27838$;
$\mathrm{Fo} = 100\times1.00486/1.28324 = 78.3$. The later \omterm{def:b1:rock-forming-minerals:groups}{olivine} is richer
in iron.

\textbf{18.} $M = 140.691$ and \qty{203.771}{g.mol^{-1}};
$V = 140.691/3.275 = 42.96$ and $203.771/4.392 =
\qty{46.40}{cm^3.mol^{-1}}$.

\textbf{19.} $M = 0.824\times140.691 + 0.176\times203.771 = 115.929 +
35.864 = \qty{151.793}{g.mol^{-1}}$; $V = 0.824\times42.96 +
0.176\times46.40 = 35.399 + 8.166 = \qty{43.565}{cm^3.mol^{-1}}$.

\textbf{20.} $\rho = 151.793/43.565 = \qty{3.484}{g.cm^{-3}}$.

\textbf{21.} $0.824\times3.275 + 0.176\times4.392 = 2.699 + 0.773 =
\qty{3.472}{g.cm^{-3}}$, \qty{0.012}{g.cm^{-3}} (\qty{0.3}{\percent}) below
the ideal-mixing value: close, as \cref{exo:b1:rock-forming-minerals:10}
predicts.

\textbf{22.} $0.12\times\qty{3484}{kg.m^{-3}}\times\qty{1}{m^3} =
\qty{418}{kg}$.

\textbf{23.} $418/2950 = \qty{14.2}{\percent}$: the dense \omterm{def:b1:rock-forming-minerals:groups}{olivine} makes up a
larger share of the mass than of the volume.

\textbf{24.} The \omterm{def:b1:rock-forming-minerals:groups}{olivine} is $\text{Fo}_{82}$: its forsterite content is
\textbf{82 per cent}.
\end{solution}
